\section{Reconhecimento Óptico de Caracteres (OCR)}
\label{sec:ocr}

\subsection{Fundamentos e Escopo}
O Reconhecimento Óptico de Caracteres (\textit{Optical Character Recognition} -- OCR) é um subcampo da Visão Computacional e da Inteligência Artificial dedicado à conversão de imagens contendo texto (impresso, manuscrito ou gravado) em dados legíveis por máquina \cite{chen2021survey}. Originalmente voltado à digitalização de documentos, o escopo foi ampliado pelo \textit{Deep Learning} para aplicações industriais e de IoT, como a leitura automatizada de medidores de utilidades públicas em ambientes não controlados \cite{jawas2018water}.

\subsection{Pipeline Tradicional}
Um sistema clássico de OCR opera em etapas interdependentes \cite{chen2021survey}:
\begin{enumerate}
    \item \textbf{Aquisição e pré-processamento}: melhoria de contraste, correção de inclinação (\textit{deskewing}), remoção de ruído e binarização; em dispositivos embarcados, etapas de segmentação semântica podem isolar a Região de Interesse (ROI) antes do reconhecimento.
    \item \textbf{Localização e segmentação de texto}: identificação das regiões de texto.
    \item \textbf{Reconhecimento de caracteres}: classificação dos segmentos ou da sequência completa.
    \item \textbf{Pós-processamento}: dicionários, correção ortográfica ou verificação sintática para mitigar erros residuais.
\end{enumerate}

\subsection{Evolução Tecnológica: dos Métodos Heurísticos ao Deep Learning}
A literatura descreve três gerações principais \cite{chen2021survey}:

\paragraph{Primeira geração (Template Matching).}
Sobreposição pixel a pixel com gabaritos estáticos; alta vulnerabilidade a variações de escala, rotação e fontes.

\paragraph{Segunda geração (Características manuais + classificadores clássicos).}
Uso de descritores como HOG e SIFT, combinados com SVM ou MLP; dependência de \textit{feature engineering}.

\paragraph{Terceira geração (End-to-End Deep Learning).}
Elimina a segmentação explícita de caracteres; modelos \textit{end-to-end} aprendem simultaneamente representações visuais e dependência contextual sequencial \cite{shi2017crnn}.

\subsection{Arquitetura CRNN e a Perda CTC}
Para sequências de comprimento variável (palavras, placas, dígitos em visores), a arquitetura padrão é a CRNN (\textit{Convolutional Recurrent Neural Network}) combinada com a perda CTC (\textit{Connectionist Temporal Classification}) \cite{shi2017crnn, graves2006ctc}.

\paragraph{Extrator convolucional (CNN).}
Camadas convolucionais extraem mapas de características bidimensionais, reduzindo a altura e preservando o eixo horizontal como sequência temporal de caracteres \cite{shi2017crnn}.

\paragraph{Modelagem de contexto.}
Tradicionalmente usam-se LSTMs/GRUs bidirecionais; em \textit{Edge AI}, porém, LSTMs impõem gargalos de memória e restrições em \textit{runtimes} C++, sendo por vezes substituídas por convoluções temporais 1D de núcleos largos que emulam o campo receptivo contextual sem o custo da recorrência \cite{shi2017crnn}.

\paragraph{CTC.}
A CTC introduz um \textit{token} oculto de espaçamento (\textit{blank}) e marginaliza todos os alinhamentos temporais válidos, permitindo treinar com sequências de tamanhos diferentes e evitando o colapso de repetições consecutivas (ex.: dígitos idênticos adjacentes) \cite{graves2006ctc}.

\subsection{Edge OCR, Quantização e TinyML}
A transição de modelos de OCR de servidores em nuvem para microcontroladores com SRAM restrita impõe o paradigma do TinyML e da Computação de Borda \cite{ademola2023resource}.

\paragraph{Quantization-Aware Training (QAT).}
O QAT simula a perda de precisão da conversão de FP32 para INT8 durante o treinamento, permitindo que os pesos da rede se adaptem e mitiguem a degradação de acurácia. Em sistemas de OCR para medidores, modelos quantizados alcançam tamanhos da ordem de centenas de kB e inferência subsegundo em hardware embarcado \cite{ademola2023resource, jawas2018water}.